\section{\sysname Design}
\label{sec:design}

The design follows the worklist from its construction and propagation across
compatible operators (\S\ref{sec:design:construct}) to its consumption by the
convolution pipeline (\S\ref{sec:design:compute}).

%% ----------------------------------------------------------------
\subsection{Worklist Construction and Reuse}
\label{sec:design:construct}

\subsubsection{Tile-Local Worklist Construction}
\label{sec:design:construct:compact}

Constructing a tight schedule already requires a scan of the output grid, but
the resulting coordinates must also retain enough spatial structure for
efficient consumption. A globally stable compaction preserves the complete
dense traversal order, but obtaining its device-wide offsets requires a prefix
scan or an additional pass. Independent appends avoid that coordination, but
scatter nearby positions across the schedule. Materializing active receptive
fields, as gather-scatter does, regularizes the compute input at the cost of
feature traffic proportional to the active count and $k^2C_{in}$.
The useful middle ground is therefore a metadata-only schedule that preserves
order locally rather than globally.

\sysname realizes this middle ground as globally unordered, locally ordered
segments. Because the number of active positions is not known before the pass,
it preallocates a coordinate buffer with capacity
$|\mathcal{G}|=H_{out}W_{out}$, the static upper bound on
$|\mathcal{W}|$, and initializes a device-resident counter for
$n_{\mathcal{W}}$. The pass then partitions the output grid into
$B_\tau\times B_\tau$ construction
tiles $\tau$. Within each tile, it applies the activity test of
Eq.~\ref{eq:mask}, writes the corresponding entries of $M_{out}$, and compacts
the active coordinates into a local sequence $\mathcal{P}_\tau$ in dense-traversal
order. When a learned gate supplies $M_{out}$ directly~\cite{dynconv,dgnet}, the
pass skips the activity test and compacts the supplied mask. With
$n_\tau=|\mathcal{P}_\tau|$, the tile reserves
$[\mathrm{base},\mathrm{base}+n_\tau)$ from the counter and copies
$\mathcal{P}_\tau$ into that interval (Algorithm~\ref{alg:construct}). After all
tiles complete, the reserved
intervals collectively store $\mathcal{W}$, and the counter records
$n_{\mathcal{W}}$. The layer shape statically determines both the buffer
capacity and the construction launch geometry, while only the device observes
the data-dependent worklist length. Construction therefore needs neither
runtime allocation nor a host-visible active count.

To relate the worklist to the useful-compute ratio of
Eq.~\ref{eq:useful-ratio}, we insert its length $n_{\mathcal{W}}$ between the
required-position count and the issued-slot count. For
$n_{\mathcal{W}}>0$,
\begin{equation}
  \eta
  = \underbrace{\frac{\|M_{out}\|_0}{n_{\mathcal{W}}}}_{\eta_{\mathrm{construct}}}
  \cdot
  \underbrace{\frac{n_{\mathcal{W}}}{|\mathcal{C}|}}_{\eta_{\mathrm{compute}}} .
  \label{eq:eta-decomp}
\end{equation}
The first factor, $\eta_{\mathrm{construct}}$, compares required positions with
worklist entries; the second, $\eta_{\mathrm{compute}}$, compares worklist
entries with issued slots. We analyze $\eta_{\mathrm{construct}}$ here and
defer $\eta_{\mathrm{compute}}$ until the compute organization is defined in
\S\ref{sec:design:compute}. The factorization decomposes $\eta$ only; all stage
costs remain reflected in throughput~$T$.

Fresh construction introduces no construction-side expansion. Its tiles
partition the output grid, and their disjoint reservations place each active
coordinate exactly once in $\mathcal{W}$. Thus
$n_{\mathcal{W}}=\|M_{out}\|_0$ and $\eta_{\mathrm{construct}}=1$ for a
nonempty worklist.
Because coverage constrains membership rather than order, the order in which
tiles reserve their intervals does not affect correctness. The within-interval
traversal order is retained for locality.

\begin{algorithm}[t]
\caption{Tile-Local Worklist Construction}
\label{alg:construct}
\begin{algorithmic}[1]
\REQUIRE Input mask $M_{in}$, output grid $\mathcal{G}$, tile size $B_\tau$
\ENSURE Mask $M_{out}$, worklist $\mathcal{W}$, length
        $n_{\mathcal{W}}=|\mathcal{W}|$
\STATE $n_{\mathcal{W}} \leftarrow 0$
\FOR{each $B_\tau \times B_\tau$ spatial tile $\tau$ \textbf{in parallel}}
  \FOR{each position $p \in \tau$ \textbf{in parallel}}
    \STATE $\mathrm{active}(p) \leftarrow \exists\, r \in \mathcal{N}(p): M_{in}[r] = 1$
    \STATE $M_{out}[p] \leftarrow \mathrm{active}(p)$
  \ENDFOR
  \STATE $\mathcal{P}_\tau \leftarrow \textsc{Compact}(\tau,\,M_{out})$
  \STATE $n_\tau \leftarrow |\mathcal{P}_\tau|$
  \STATE \COMMENT{Atomically reserve $n_\tau$ worklist entries.}
  \STATE $\mathrm{base} \leftarrow n_{\mathcal{W}}$
  \STATE $n_{\mathcal{W}} \leftarrow n_{\mathcal{W}} + n_\tau$
  \FOR{each $0\le j<n_\tau$ \textbf{in parallel}}
    \STATE $\mathcal{W}[\mathrm{base}+j] \leftarrow \mathcal{P}_\tau[j]$
  \ENDFOR
\ENDFOR
\end{algorithmic}
\end{algorithm}

\sysname deliberately forgoes a global worklist order, allowing construction
tiles to reserve contiguous segments independently without device-wide ordering
coordination. Within each segment, active coordinates follow the tile's dense
traversal, giving compute a spatially grouped coordinate run but offering no
locality guarantee across segment boundaries.
As shown in \S\ref{sec:eval:efficiency}, tile-local ordering achieves
compute-stage throughput comparable to that under global ordering; avoiding the
additional ordering pass is therefore the better end-to-end trade-off.

\subsubsection{Cross-Operator Worklist Reuse}
\label{sec:design:reuse}

Construction has a channel-independent grid-scan cost. Each pass examines all
$H_{out}W_{out}$ positions and, when deriving $M_{out}$ from $M_{in}$, performs
$O(k^2)$ mask work per position, whereas required convolution work scales as
$O(\alpha H_{out}W_{out}k^2C_{in}C_{out})$. Repeating this scan across
consecutive same-grid layers that introduce no new required positions is
therefore redundant. \sysname instead reuses one constructed worklist across
each compatible layer sequence, amortizing its construction cost over the
sequence.

For layer $l$, let $\mathcal{G}^l$ denote its output grid, $M^l_{out}$ its
output mask, $\mathcal{W}^l$ its worklist, and $n^l_{\mathcal{W}}$ its
device-resident length. \sysname propagates the mask and worklist metadata
alongside the feature tensor. From layer semantics, \sysname statically infers
that consecutive layers are reuse-compatible when
$\mathcal{G}^{l+1}=\mathcal{G}^l$ and
\begin{equation}
  \{p\in\mathcal{G}^{l+1} \mid M^{l+1}_{out}[p]=1\}
  \subseteq
  \{p\in\mathcal{G}^l \mid M^l_{out}[p]=1\} .
  \label{eq:reuse-condition}
\end{equation}
Because $\mathcal{W}^l$ already satisfies coverage for $M^l_{out}$,
Eq.~\ref{eq:reuse-condition} guarantees coverage for $M^{l+1}_{out}$ as well.
\sysname therefore sets $\mathcal{W}^{l+1}=\mathcal{W}^l$ and
$n^{l+1}_{\mathcal{W}}=n^l_{\mathcal{W}}$ without reconstruction. \sysname
applies the same condition recursively to subsequent layers, forwarding the
original $\mathcal{W}^l$ and $n^l_{\mathcal{W}}$ through the entire compatible
sequence.

Equality in Eq.~\ref{eq:reuse-condition} covers same-grid $1\times1$
convolutions, mask-preserving activations, and inference-time normalization and
leaves $\eta_{\mathrm{construct}}$ unchanged. Strict inclusion results from
mask-narrowing operators in the upstream policy, such as the activation-fused
update truncation that DeltaCNN uses to increase downstream
sparsity~\cite{deltacnn}. \sysname therefore retains $\mathcal{W}^{l+1}$ as a
valid superset, trading the resulting lower $\eta_{\mathrm{construct}}$ for one
fewer construction pass without violating coverage.

A change in the output grid violates the geometry constraint, while
receptive-field expansion may violate Eq.~\ref{eq:reuse-condition} by
introducing positions outside the upstream mask; either case invokes the
construction pass of \S\ref{sec:design:construct:compact} to emit a new tight
worklist. A lower $\eta_{\mathrm{construct}}$ alone does not trigger
reconstruction because tightening a valid superset would pay another grid scan.
Reuse avoids repeated construction inside each compatible sequence, but common
receptive-field-expanding layers such as $3\times3$ convolutions create rebuild
points throughout the network. The resulting device-resident worklist lengths
vary across layers and frames, so compute still faces a dynamic problem size.

%% ----------------------------------------------------------------
\subsection{Variable-Length Worklist Execution}
\label{sec:design:compute}

Construction and reuse fix $\eta_{\mathrm{construct}}$ before compute begins;
the implicit-GEMM decomposition controls the remaining
$\eta_{\mathrm{compute}}$ and also affects throughput~$T$. For a standard
ungrouped batch-1 convolution, dense implicit GEMM uses
$M=H_{out}W_{out}$ output-position rows, $N=C_{out}$ output-channel columns, and
reduction depth $K=k^2C_{in}$. It decomposes the output into
$B_M\times B_N$ tiles and traverses each tile's reduction in $B_K$-wide steps.
\sysname preserves the $N$ and $K$ dimensions and this tiled pipeline, but
changes the $M$-dimension coordinate source: instead of enumerating the dense
output grid, it reads the $n_{\mathcal{W}}$ coordinates in $\mathcal{W}$.
Compute therefore groups $\mathcal{W}$ into $B_M$-coordinate tiles, padding
only the final tile. For $n_{\mathcal{W}}>0$, this issues
$|\mathcal{C}|=B_M\lceil n_{\mathcal{W}}/B_M\rceil$ position slots, giving
\begin{equation}
  \eta_{\mathrm{compute}}
  = \frac{n_{\mathcal{W}}}
         {B_M\left\lceil n_{\mathcal{W}}/B_M\right\rceil} .
  \label{eq:eta-compute}
\end{equation}
The final tile contains at most $B_M-1$ padding slots, making this overhead most
pronounced for a short worklist. The same $B_M$ also controls data reuse, block
efficiency, and exposed parallelism, so maximizing $\eta_{\mathrm{compute}}$
need not maximize throughput~$T$.

For dense convolution, $(M,N,K)$ follow from the layer shape, allowing the
backend to select $(B_M,B_N,B_K)$ before execution. In the worklist path,
the $M$-dimension extent is $n_{\mathcal{W}}$, which varies across invocations
and remains unknown to the host. This creates two problems: executing the
current worklist without exposing its length to the host, and selecting a
decomposition suited to that length.
\S\ref{sec:design:compute:persistent} addresses the former with persistent
execution, while \S\ref{sec:design:adaptive} addresses the latter through
trace-calibrated selection.

\subsubsection{Persistent Worklist-Driven Convolution}
\label{sec:design:compute:persistent}

Concretely, the dense path maps row $m$ to the $m$-th coordinate $p_m$ of the
contiguous traversal
($0\le m<|\mathcal{G}|$); the worklist path instead reads $\mathcal{W}[m]$
($0\le m<n_{\mathcal{W}}$). Each block loads a $B_M$-slot coordinate group,
decodes the receptive-field bases of its non-$\bot$ entries, and computes their
$B_N$ output channels. The coordinates determine the input gathers and output
stores, while the kernel retains the regular $B_K$-tiled reduction over~$K$.
It therefore reads directly from the dense feature tensor without materializing
receptive-field patches. Processing $\mathcal{W}$ in contiguous ranges carries
construction's within-segment grouping into these gathers.

This substitution defines the scheduled coordinates, but their number remains
unknown to the host. \sysname executes the device-resident worklist with an
occupancy-sized persistent grid. At initialization, it uses each kernel
configuration's concurrent residency to fix the corresponding block count.
Selecting a configuration therefore selects a grid bounded by device capacity
rather than $n_{\mathcal{W}}$. Once launched, the blocks read the current
device-resident length and derive the base work-item count
\begin{equation}
  n_{\mathrm{base}} = \left\lceil \frac{n_{\mathcal{W}}}{B_M} \right\rceil
  \left\lceil \frac{C_{out}}{B_N} \right\rceil ,
  \label{eq:workitems}
\end{equation}
and cover all position and channel tiles with a grid-stride loop
(Algorithm~\ref{alg:compute}). The device-resident length determines the loop
bound and final-tile padding, while the host launch remains independent of that
length. The count $n_{\mathrm{base}}$ includes channel tiles because it counts
thread-block work items; $|\mathcal{C}|$ instead counts each issued position
once and therefore has no channel-tile factor.

Sizing the grid from $n_{\mathcal{W}}$ on the host would synchronize at every
operator. Sizing it for the full-grid upper bound avoids that readback, but
submits blocks in proportion to the dense grid even when the worklist is short.
\sysname instead bounds submitted blocks by concurrent device residency and
resolves the actual work only after launch. Persistent execution thereby
handles any device-resident row extent; choosing the configuration that handles
it efficiently remains a separate problem.

\begin{algorithm}[t]
\caption{Persistent Worklist-Driven Implicit GEMM}
\label{alg:compute}
\begin{algorithmic}[1]
\REQUIRE Device-resident worklist $\mathcal{W}$ and length
         $n_{\mathcal{W}}=|\mathcal{W}|$
\REQUIRE Input $\mathbf{x}$ and weights $\mathbf{w}$
\ENSURE Output $\mathbf{y}$ (updated at worklist coordinates)
\STATE $n_{\mathcal{W}} \leftarrow$ \textbf{read} device-resident length
\STATE \COMMENT{Dense implicit GEMM uses host-known $M=H_{out}W_{out}$.}
\STATE $n_{\mathrm{base}} \leftarrow \lceil n_{\mathcal{W}}/B_M \rceil
       \lceil C_{out}/B_N \rceil$
\FOR{$\mathrm{tile}=\mathrm{block\_id},\;
      \mathrm{block\_id}+\mathrm{grid\_size},\;\ldots<n_{\mathrm{base}}$}
  \STATE $(b_m,b_n) \leftarrow \mathrm{Decompose}(\mathrm{tile})$
  \FOR{each row $0\le j<B_M$}
    \STATE $p_j \leftarrow \mathcal{W}[b_mB_M+j]$ if $b_mB_M+j<n_{\mathcal{W}}$
    \STATE $p_j \leftarrow \bot$ otherwise
  \ENDFOR
  \STATE \COMMENT{Dense sets $p_j \leftarrow p_{b_mB_M+j}$ and bounds against $M$.}
  \STATE Decode each non-$\bot$ $p_j$ to its receptive-field base address
  \STATE Initialize accumulator $\mathbf{Acc} \leftarrow \mathbf{0}_{B_M \times B_N}$
  \FOR{each $B_K$-wide reduction step over $K = k^2 C_{in}$}
    \STATE $\mathbf{x}_{\mathrm{tile}} \leftarrow$ gather for non-$\bot$ rows
    \STATE $\mathbf{w}_{\mathrm{tile}} \leftarrow$ load corresponding weight slice
    \STATE $\mathbf{Acc} \leftarrow \mathbf{Acc} + \mathbf{x}_{\mathrm{tile}} \cdot \mathbf{w}_{\mathrm{tile}}$
  \ENDFOR
  \STATE Store $\mathbf{Acc}$ to $\mathbf{y}$ at the non-$\bot$ coordinates $p_j$
\ENDFOR
\end{algorithmic}
\end{algorithm}

\subsubsection{Trace-Calibrated Adaptive Decomposition}
\label{sec:design:adaptive}

Persistent execution handles every worklist length, but the fastest candidate
remains length-dependent because the decomposition jointly controls
$\eta_{\mathrm{compute}}$ and throughput~$T$. For operator $i$ with output
grid $\mathcal{G}_i$ and worklist length $n^i_{\mathcal{W}}$, its active ratio
and normalized worklist length are
\begin{equation}
  \alpha_i = \frac{\|M^i_{out}\|_0}{|\mathcal{G}_i|}, \qquad
  \lambda_i = \frac{n^i_{\mathcal{W}}}{|\mathcal{G}_i|} .
  \label{eq:lambda}
\end{equation}
Fresh construction gives $\lambda_i=\alpha_i$, whereas superset reuse gives
$\lambda_i\ge\alpha_i$. Because a worklist candidate processes the scheduled
length rather than only the active set, \sysname discretizes $\lambda_i$ into
levels $\ell$, each represented by an absolute length $n_{i,\ell}$.
Let $\mathcal{Q}_i$ be the candidates enumerated for operator $i$ on the
deployment device, and let $\widehat{t}_i(q;n)$ be the measured latency of
candidate $q$ at length $n$. The worklist candidates in $\mathcal{Q}_i$ span
$(B_M,B_N,B_K,S_k)$, where $S_k$ partitions the reduction to expose additional
parallelism. For a known level, the best candidate follows from standard
autotuning:
\begin{equation}
  q^*_{i,\ell} = \arg\min_{q\in\mathcal{Q}_i}
  \widehat{t}_i(q;\,n_{i,\ell}) .
  \label{eq:calib}
\end{equation}

The candidate knobs affect $\eta_{\mathrm{compute}}$ and $T$ differently.
$B_M$ trades final-tile padding and the number of independent position tiles
against intra-tile reuse and per-block efficiency. $B_N$ and $B_K$ tune the
regular channel and reduction datapath without changing $|\mathcal{C}|$. When
position tiles alone expose too little parallelism, $S_k$ partitions $K$ and
multiplies the work-item count in Eq.~\ref{eq:workitems}, also without changing
$|\mathcal{C}|$. This split can raise throughput on short worklists, but rereads
coordinates and requires output zeroing and atomic accumulation. Calibration
therefore measures every candidate-specific cost on worklists with the
tile-local structure produced by construction, excluding only the construction
pass shared by all candidates.

When full-grid execution preserves the operator semantics, $\mathcal{Q}_i$
also includes a dense candidate in the same measured-latency objective; it
forwards the mask and worklist metadata unchanged for downstream operators.

Calibration therefore resolves $q^*_{i,\ell}$ offline. The remaining challenge
is selecting $\ell$ at runtime without reading the current
$n_{\mathcal{W}}$ back at every operator.

\sysname avoids the readback by separating the current worklist bound from an
approximate configuration signal. Persistent compute reads
$n_{\mathcal{W}}$ on the device, so the host needs only a predictor of the
corresponding level, not the exact bound. Event and temporal workloads provide
one when two properties hold:
their layer masks derive from a shared network-input mask through
receptive-field propagation, and successive input masks correlate in time.
Under these conditions, a completed observation of global input activity from
an earlier frame provides a predictor of current per-operator lengths for
decomposition selection.

Let $M^f_0$ denote this network-input mask for frame $f$ on grid
$\mathcal{G}_0$, and let
\begin{equation*}
  \beta_f = \frac{\|M^f_0\|_0}{|\mathcal{G}_0|}, \qquad
  b_f = \mathrm{Bucket}(\beta_f).
\end{equation*}
\sysname learns a per-operator mapping $\mathcal{M}_i$ from input bucket to
worklist-length level. Combined with the candidate table
$\mathcal{L}_i[\ell]=q^*_{i,\ell}$, a completed observation selects
\begin{equation}
  q_i(f) = \mathcal{L}_i\!\left[\mathcal{M}_i(\hat{b}_f)\right] ,
  \label{eq:select}
\end{equation}
where $\hat{b}_f$ is the bucket derived from the most recent active-count
transfer completed by the start of frame $f$. The observation may be one or
more frames old, but it affects only the chosen configuration; the current
device-resident $n^i_{\mathcal{W}}$ still bounds the worklist execution.

A second offline pass learns $\mathcal{M}_i$ from representative temporal
traces. It first fixes the deployment reuse decisions, then records each
frame's input bucket $b_f$ together with the worklist length
$n^i_{\mathcal{W}}(f)$ presented to every operator. Recording this schedule
length, rather than only the cardinality of
$M^i_{out}$, incorporates the superset schedules created by reuse. For each
bucket, \sysname maps the
median normalized length to a level:
\begin{equation}
  \mathcal{M}_i(b) = \mathrm{Level}\!\left(
    \operatorname*{median}_{f\,:\,b_f = b} \lambda_i(f) \right) .
  \label{eq:mapfit}
\end{equation}
The median targets the representative length rather than allowing a few
outlying frames to determine the configuration. Each operator needs its own
mapping because grid resolution, receptive-field size, and propagation depth
transform the same input activity differently. The two calibration passes now
compose directly: Eq.~\ref{eq:mapfit} predicts a worklist-length level from
input activity, and Eq.~\ref{eq:calib} supplies the fastest measured candidate
for that level.

The online loop applies this mapping using only completed observations
(Algorithm~\ref{alg:adapt}). Each operator starts from a calibrated default. At
the beginning of a frame, the runtime updates configurations from the most
recent completed transfer, or retains them if none is available. It then
generates $M^f_0$, enqueues an asynchronous transfer of its active count, and
launches construction and compute. An enqueued observation becomes eligible
only in a later frame; if the observation buffer is full, the runtime drops the
new transfer. No step therefore waits for an activity observation.

\begin{algorithm}[t]
\caption{Trace-Calibrated Nonblocking Configuration Selection}
\label{alg:adapt}
\begin{algorithmic}[1]
\REQUIRE Candidate tables $\mathcal{L}_i$ and representative traces
\ENSURE Per-operator candidates $q_i(f)$ without host synchronization
\STATE \textbf{Offline trace mapping:}
\FOR{each frame $f$ of the traces}
  \STATE $b_f \leftarrow \mathrm{Bucket}(\|M^f_0\|_0 / |\mathcal{G}_0|)$
  \FOR{each operator $i$}
    \STATE Append $n^i_{\mathcal{W}}(f)/|\mathcal{G}_i|$ to $R_i[b_f]$
  \ENDFOR
\ENDFOR
\STATE $\mathcal{M}_i(b) \leftarrow \mathrm{Level}(\mathrm{median}\, R_i[b])$ for all $i,b$
\STATE \textbf{Online (per sequence):}
\STATE $q_i \leftarrow$ calibrated default entry of $\mathcal{L}_i$ for all $i$
\FOR{each frame $f$}
  \IF{an observation has completed}
    \STATE $\hat{b}_f \leftarrow$ bucket of the most recent completed observation
    \STATE $q_i \leftarrow \mathcal{L}_i[\mathcal{M}_i(\hat{b}_f)]$ for all $i$
  \ELSE
    \STATE Retain current $q_i$ \COMMENT{do not wait for a pending transfer}
  \ENDIF
  \STATE Generate $M^f_0$ and its device-resident count
  \STATE Enqueue asynchronous count transfer; drop if the buffer is full
  \STATE Run each $q_i$; worklist candidates use current $n^i_{\mathcal{W}}$
\ENDFOR
\end{algorithmic}
\end{algorithm}

Because every candidate satisfies coverage and each worklist candidate still
reads the current $n^i_{\mathcal{W}}$, a stale, missing, or inaccurate
observation can affect latency but not coverage.

Learned-gating models do not satisfy the shared-input-mask condition needed by
the feedback mapping. Their network input remains dense, while each gated block
derives its own mask from intermediate features, so global input activity
conveys no information about per-operator worklist lengths. \sysname assigns
each such operator a fixed level calibrated from its gate statistics because
tracking every gate would require one observation per operator rather than one
per frame. Temporal workloads thus select calibrated candidates from completed
observations, while learned-gating workloads retain fixed levels. In both
cases, the current $n_{\mathcal{W}}$ bounds every worklist execution without
synchronizing the per-frame path.
